% Complete port of PROMOTION_WIDTH.qmd; all auxiliary symbols are local.
% Book readout W^{(L+1)} is the compact paper's w.
\section{Polynomial sufficient widths for Legendre and Taylor}
\label{cp:poly:width-section}

The sufficient widths for Legendre and finite-panel Taylor can be made
polynomial in the exposed problem parameters while preserving their
storage bounds. This requires a finite source-probability argument and a
stronger Legendre approximation estimate. The logarithmic relative-error
estimate alone would not give polynomial dependence on a requested error.

Use the model, initialization, activations, label condition and trajectory
norm in \hyperref[cp:setup]{the common setup}. All auxiliary notation in this technical
section is local to it. Put

\[
X=\beta^L,\qquad \lambda=\gamma/m\in(0,1],\qquad
\chi=1+\lambda^{-1},\qquad S=16Y/\lambda,\qquad \ell_n=\log(en).
\]

Thus $X\ge100$, $Y/\lambda\le X^{-30}$ and $S\le X^{-26}$.
The scalar $S$ is a normalization, not additional state. For $0<\rho<1/4$
and a declared list of $N\ge m$ inputs define

\[
u_\rho=\max\left\{1,\left\lceil
 \frac{\log(4emL/\rho)}{\log(64e^2)}\right\rceil\right\},\qquad
A_{N,\rho}=X^{4000}\chi^8(N+d+u_\rho+1)^8,
\]

\begin{equation}
N_{\rm src}(N,\rho)
=\left\lceil(A_{N,\rho}/\rho)^{40000}\right\rceil.
\label{cp:poly:eq:source-width}
\end{equation}

The large common exponent is convenient for a uniform finite ledger; it
does not claim practical sharpness. In particular this gate has no $Y^{-1}$.

\begin{theorem}[Polynomial onset with the original storage powers]
\label{cp:poly:width}

For $0<\delta,\varepsilon<1$, a sufficient width for

\[
\Pr\left\{
 \|f_{\rm model}-f_n\|_{\mathcal X}
 \le\varepsilon\|f_n-\widetilde f_n\|_{\mathcal X}
 \right\}\ge1-\delta
\]

is bounded by

\[
\left[c_{\rm wid}\beta^L
 (1+m+d+p+\gamma^{-1}+Y^{-1}+\delta^{-1}+\varepsilon^{-1})
 \right]^{c_{\rm wid}}
\]

with one absolute finite numerical constant $c_{\rm wid}$, local to this theorem. For Legendre take $p=0$ and omit
$Y^{-1}$; its domain is the whole input sphere. For Taylor, $p$ is the
number of passive inputs declared before initialization and its domain is
that panel together with the training inputs. The Legendre and Taylor
storage bounds in \eqref{cp:legendre-size} and
\eqref{cp:panel-size} hold with activation/depth coefficients bounded
by fixed powers of $X$. Taylor also has absolute error $Y/n$.

The constructions, physical optimizer, initialization-only provenance and
label condition are unchanged. Confidence applies to each individual width.
The inverse-label gate for Taylor is needed only to simplify its exact
source degree to a multiple of $\log(en)$. This result makes no polynomial
onset claim for Harmonic, for the nonlinear feature-learning certificate,
or for the stronger coefficient in \eqref{cp:eq-dense-lower}, and gives no
pointwise-in-time or endpoint denominator bound.

\end{theorem}

We prove a finite source lemma, obtain a finite lower bound for actual
dense variability, improve the Legendre history defect from order $q^{-2}$
to $q^{-3}$, and then assemble the two approximation estimates. The exact
gates at the end make the polynomial assertion explicit.

\subsection{Finite initialization and analytic source probability}
\label{cp:poly:source-section}

The coefficient recurrences in \eqref{cp:src-basic},
\eqref{cp:src-absorption-coefficients},
\eqref{cp:src-allowance}, \eqref{cp:src-response-coefficients},
\eqref{cp:src-UV} and \eqref{cp:src-K} are used with precisely
their displayed definitions. In this subsection their symbols
$H_j,P_j,f_j,k_j,\tau_j,A_*,D_*,\eta_*,\mathcal B,\mathcal K,
D_W,K_{\rm src},U_{\rm fin}$ have those meanings. We quantify the
finite-input proof; no angular source or whole-sphere complex domain is used.
The scalar envelopes used in those recurrences are, locally,
\[
b=\max_j|\phi_j(0)|,\quad
s=\max\{1,\sup_{j,|\operatorname{Im}z|\le a/2}|\phi'_j(z)|\},\quad
t_2=\max\{1,\sup_{j,|\operatorname{Im}z|\le a/2}|\phi''_j(z)|\}.
\]
The additional recurrence coefficients, including $C_{\rm abs}$ and
$W_G$, always mean the displayed quantities in
\eqref{cp:src-absorption-coefficients}; they remain local to this argument.
The pre-gated carriers are $k_a^{(L)}=w$ and
$k_a^{(j)}=W^{(j+1)\top}\delta_a^{(j+1)}$ for $j<L$.
In the stopped-flow estimates below, $\rho(t)=\|r(t)\|_2/\sqrt m$
denotes residual RMS; the constant $\rho$ in the width and probability
gates denotes the failure allowance.
The Euclidean mobility coordinates used for real parameter differences are
$\theta=(W^{(1)}/\sqrt n,W^{(2)},\ldots,W^{(L)},w/\sqrt n)$,
with the ordinary Euclidean/Frobenius block norm.

\begin{lemma}[Finite training and panel sources]
\label{cp:poly:source}

At $n\ge N_{\rm src}(N,\rho)$, with probability at least $1-\rho$,
the ambient real fitting event holds and the following source conclusions
hold. The four families are exactly \eqref{cp:source-families};
forward inputs range over the $N$ declared inputs and backward inputs over
the $m$ training inputs.

For the training list alone one may use the closed rectangle

\begin{equation}
[-r,T+r]+i[-r,r],\qquad
T=32\ell_n/\lambda,\qquad
r=[u_\rho X^{100}(1+\lambda)\sqrt{\ell_n}]^{-1}.
\label{cp:poly:eq:training-radius}
\end{equation}

All training sources and the responses
$R_{ba}^{(j)}=D_\Theta z_b^{(j)}\nabla_\Theta(nf_a)$, where

\[
\Theta=(W^{(1)},\sqrt nW^{(2)},\ldots,\sqrt nW^{(L)},w),
\]

are holomorphic near this rectangle. The source RMS bounds are
\eqref{cp:source-rms}, their coordinates are bounded by
$X^6\sqrt n$, and $|f_n|\le16X^6Y/\lambda$. Moreover

\[
\max_{a,b,j,i}|R_{ba,i}^{(j)}|\le SU_{\rm fin}\sqrt{\ell_n}
\]

on that rectangle. On either complex source domain the first-layer
operator divided by $\sqrt n$ and all hidden operators have norm at most
ten, the complex training residual RMS is at most $2Y$, and every declared
preactivation has imaginary part below $3a/8$. The training coordinate
maxima are bounded by $K_{\rm src}\sqrt{\ell_n}$ for preactivations and
$SK_{\rm src}\sqrt{\ell_n}$ for carriers. These strict strip and
coordinate bounds are included in the event, not inferred from an RMS
bound. On the whole real training trajectory, including its
limit,

\begin{equation}
\max_{a,j,i}|k_{a,i}^{(j)}|\le32X^{21}(Y/\lambda)\sqrt{\ell_n},
\qquad
\max_{a,b,j,i}|R_{ba,i}^{(j)}|
 \le YX^{110}\chi^3(1+\sqrt{\ell_n}).
\label{cp:poly:eq:real-response}
\end{equation}

For the $N$-input panel put $a_0=\min(a,1)$ and

\begin{equation}
b_n=\log(e+\ell_n),\quad B_n=1+(S^2/\eta_*)b_n,\quad
r_n=\min\left\{\frac{a_0}{64YSU_{\rm fin}\sqrt{\ell_n}},
 \frac1{X^{300}(1+\lambda B_n)\sqrt{b_n}}\right\}.
\label{cp:poly:eq:panel-radius}
\end{equation}

There is a deterministic disk cover and half-radius panel partition of
$[0,T]$, with the same source envelopes, whose number of panels obeys

\begin{equation}
J\le1+X^{400}\left[(Y/\lambda)^2\sqrt{\ell_n}
                      +\lambda^{-1}\log(e+\ell_n)\right].
\label{cp:poly:eq:panel-count}
\end{equation}

Neither this radius nor its panel-count coefficient depends on confidence.

\end{lemma}

\begin{proof}

First quantify initialization. Let
$H=\max_{0\le j\le L}\sqrt{Q^{(j)}_{11}}\ge1$ and let $b,s$
have their meanings in the source proof. Define

\[
R_0=2H,\quad K_0=b+sR_0,\quad D_0^{\rm cov}=2sb+4s^2R_0,
\quad T_1=\sum_{i=0}^{L-1}s^i,\quad T_2=\sum_{i=0}^{L-1}s^{2i},
\]

\[
C_0^{\rm cov}=T_2(2K_0+1+D_0^{\rm cov}T_1),\quad
e_0=\min\{(8T_1)^{-1},\gamma/(4mC_0^{\rm cov})\},\quad
h_0=\min\{1/2,H/[4(8s)^L]\}.
\]

The Gaussian Lipschitz inequality proved in
\hyperref[cp:src-Gaussian]{the Gaussian proof} gives, for an $a$-Lipschitz nonnegative
Gaussian function $F$,
$\operatorname{Var}F\le a^2$ and
$0\le(\mathbb EF^2)^{1/2}-\mathbb EF\le a$.
Consequently, for $e\ge2a$,

\[
\Pr\{|F-(\mathbb EF^2)^{1/2}|>e\}\le2e^{-e^2/(8a^2)}.
\]

For iid Gaussian row pairs with marginal standard deviations at most $R_0$,
the empirical RMS of $\phi(Z_a)$ is $sR_0/\sqrt n$-Lipschitz.
The RMS of $\phi(Z_a)\pm\phi(Z_b)$ is
$2sR_0/\sqrt n$-Lipschitz, including singular covariance pairs.
Polarization therefore gives covariance-entry error at most
$(2K_0+1)e_0$ if both such norms have error at most $e_0\le1$.

To propagate covariance errors without dividing by a variance, consider
two Gaussian pairs with standard deviations $\sigma_a,\sigma_b$ and
$\sigma'_a,\sigma'_b$, all at most $R_0$, and set
$\Delta=\max_i|\sigma_i-\sigma'_i|$.
Shrink both pairs to marginal deviations
$\tau_i=\min(\sigma_i,\sigma'_i)$, preserving their correlations.
Product subtraction and Cauchy--Schwarz bound the total scaling cost by
$2sK_0\Delta$. At these common marginals Gaussian integration by parts
gives the correlation derivative
$\tau_a\tau_b\mathbb E[\phi'(Z_a)\phi'(Z_b)]$.
The formula follows by differentiating the nonsingular density and
integrating its mixed second derivative; linear growth and bounded
derivatives justify the integrations. Continuity covers correlations
$\pm1$ and zero marginal variances. Since

\[
\tau_a\tau_b|\rho_c-\rho'_c|
 \le|C_{ab}-C'_{ab}|+2R_0\Delta,
\]

the transformed entry error is at most
$s^2|C_{ab}-C'_{ab}|+D_0^{\rm cov}\Delta$.
The displayed inequality follows by writing the correlation difference
over the larger of $\sigma_a\sigma_b,\sigma'_a\sigma'_b$ and using
$|C'_{ab}|\le\sigma'_a\sigma'_b$.

Condition at each layer on the preceding initialized layers. On a sphere
$h_0$-net of size at most $(1+2/h_0)^d$, and the training list, let
$\Delta_j$ denote the largest RMS error and $E_j$ the largest training
Gram-entry error. The preceding estimates give

\[
\Delta_j\le e_0+s\Delta_{j-1},\qquad
E_j\le(2K_0+1)e_0+s^2E_{j-1}+D_0^{\rm cov}\Delta_{j-1}.
\]

Thus $\Delta_j\le e_0T_1\le1/8$ and
$E_j\le e_0C_0^{\rm cov}\le\gamma/(4m)$. This also keeps all
conditional standard deviations below $R_0$.
With

\[
\Xi=\log\frac{16L(m+1)^2}{\rho_0}
                 +d\log(1+2/h_0),
\]

a sufficient initialized fitting width at failure $\rho_0$ is

\begin{equation}
\left\lceil\max\left\{1,
 \frac{\log(8L/\rho_0)}{8-2\log9},
 \frac{d\log9+\log(8/\rho_0)}{8-\log9},
 \frac{32s^2R_0^2\Xi}{e_0^2}\right\}\right\rceil.
\label{cp:poly:eq:initialization}
\end{equation}

Indeed the layer union has at most $3m^2+(1+2/h_0)^d$ norm tests per
layer and costs at most $\rho_0/2$. The operator-net tails in
\hyperref[cp:fit]{the fitting proof} cost the other half. The operator event makes the
initialized layer-$j$ feature map in neuron RMS $(8s)^j$-Lipschitz;
its off-net increment is at most $H/4$. This proves
\eqref{cp:eq-initialization-event}. The deterministic fitting proof then
applies with the existing label cap. Unrolling the displayed coefficients
gives

\[
R_0,K_0\le X^3,\quad D_0^{\rm cov}\le X^4,\quad
T_1\le X^2,\ T_2\le X^3,\ C_0^{\rm cov}\le X^{10},\quad
e_0^{-1}\le X^{11}\chi.
\]

Since $32s^2R_0^2\le X^8$ and
$\log(1+2/h_0)\le3L\log\beta$, the width is at most
$\lceil X^{32}\chi^2[dL\log\beta+\log(16L(m+1)^2/\rho_0)]\rceil$.
This proof imposes no ordering between $m$ and $d$.

For source initialization choose $\rho_0=\rho/64$.
Conditional Gaussian tails over the declared list give
$Z_0=2H\sqrt{2\log(64nNL/\rho)}$ as an initial preactivation envelope.
Deleting at most $u=u_\rho$ neurons changes each subsequent initialized
feature vector in Euclidean norm by at most
$D_{\rm init}=(8s)^L\sqrt u(b+sZ_0)$, by zero embedding and forward
subtraction. The full training Gram starts above $3\gamma/4$ and the
source gate gives $D_{\rm init}/\sqrt n<\sqrt\lambda/8$.
Thus every cavity's normalized least singular value remains above
$(\sqrt3/2-1/8)\sqrt\lambda>\sqrt{\lambda/2}$, preserving a strict
initial fitting margin. Pairing a projected lower difference with a fresh row gives tail
$2\exp[-n^{4/5}/(2D_{\rm init}^2)]$ at threshold $n^{-1/10}$;
the projection is made before conditioning on that fresh row.
There are at most $uLn^uNLn$ such tests. Initially each exponential
budget has conditional mean below four and second moment below fifteen:
$\mathbb E e^{\eta_*|Z|}\le2e^{1/2}$ and
$\mathbb E e^{2\eta_*|Z|}\le2e^2$, since $\eta_*2H\le1$.
Conditional Chebyshev costs at most $15mL/(16n)$ for the layer averages
to remain below eight. Together these initialization failures are below
$\rho/4$ at \eqref{cp:poly:eq:source-width}. Cavities have their
own initial tests; failed tests define zero reference coefficient paths.

We next supply the quantitative map coefficients that replace qualitative
polynomials in the source proof. Use its independently stopped normalized
equation \eqref{cp:src-retained} with
$\bar u=(\Theta-\Theta_0)/S$ and $\bar F_a=nf_a/S$.
Zero readout gives $\bar F_a=\bar u_{w}^\top h_a^{(L)}$;
there is no inverse $S$ in its derivatives. At a cavity, freeze scalar
forward and scaled reverse controls and let $g\sim N(0,I_r/n)$,
$r\le2un$, concatenate the omitted interior roots. First-layer incoming
rows are represented only in their own uniformly tested scalar controls;
there is no retained reverse port below that layer.

Write $\kappa_0=1/4000$, $\Gamma=\sqrt{\mathcal K}$,
$\tau_*=\max_j\tau_j$, $f_*=\max_jf_j$, and
$T_0=64/\lambda$. The propagator coefficient is

\[
J_0=2e^{1/4}(2\mathcal B)^{1/4000},\qquad
\|J(t,s)\|\le J_0n^{\kappa_0}.
\]

Define, in the order displayed,

\[
M_0=\eta_*^{-1}[1+\log(2\mathcal B)],\quad
A_0=b+sM_0,\quad B_0=sM_0,\quad E_0=A_*+SD_*M_0,
\]

\[
\begin{gathered}
T_L^0=2sH_L+2t_2M_0V_L^0,\quad
T_j^0=10sT_{j+1}^0+2s\tau_{j+1}^2H_j+2t_2M_0V_j^0,\\
D_c=\max\{\lambda s\max_jV_j^0/4,\lambda\max_jT_j^0/8\}.
\end{gathered}
\]

The source budget bounds every scaled carrier by $M_0\ell_n$;
the controls have amplitudes $A_0\ell_n,B_0\ell_n$ and speed at most
$D_c\sqrt n\ell_n$. The first-variation coefficients are

\[
v_0=J_0\{A_0(2T_0\Gamma\tau_*+E_0)+Sf_*B_0\},\quad
z_j^0=P_j(Sv_0+A_0),\quad h_j^0=sz_j^0,
\]

\begin{equation}
k_L^0=v_0,\quad d_j^0=sk_j^0+t_2M_0z_j^0,\quad
k_j^0=10d_{j+1}^0+S\tau_{j+1}v_0+B_0\quad(j<L).
\label{cp:poly:eq:map-coefficients}
\end{equation}

They bound the parameter, forward and backward linear maps from $g$ by,
respectively,
$\sqrt u n^{\kappa_0}v_0\ell_n^2$,
$\sqrt u n^{\kappa_0}(z_j^0,h_j^0)\ell_n^2$ and
$\sqrt u n^{\kappa_0}(k_j^0,d_j^0)\ell_n^3$.
To verify the parameter coefficient, the three differentiated port forces
are

\[
-\frac2{mn}\nabla_{\bar u}\bar F_a\bar\delta_a^{(j+1)\top},\quad
-\frac2m\bar r_a(D_{\bar u}\bar\delta_a^{(j+1)})^\top,\quad
-\frac2m\bar r_a(D_{\bar u}h_a^{(j-1)})^\top.
\]

Their norms use $\Gamma,\tau_*,E_0\ell_n,Sf_*$, and Duhamel integration
uses $2\int\rho/S\le1$ and contour length $T_0\ell_n$.
The differentiated forward/backward identities in
\hyperref[cp:src-local-lemma]{the local-map proof} then give exactly the remaining
recurrences. A reference carrier enters an additive forcing once;
homogeneous propagation is $10s$.

For a control difference replace $A_0,B_0$ by one in the forcing, keeping
all reference coefficients fixed:

\[
v_c=J_0(2T_0\Gamma\tau_*+E_0+Sf_*),\quad
z_j^c=P_j(Sv_c+1),\quad h_j^c=sz_j^c,
\]

\[
k_L^c=v_c,\quad d_j^c=sk_j^c+t_2M_0z_j^c,\quad
k_j^c=10d_{j+1}^c+S\tau_{j+1}v_c+1.
\]

The maximum $C_c$ of these coefficients bounds the map difference by
$C_c\sqrt u n^{\kappa_0}\ell_n^2$ times the supremum control error.
Linearity in controls proves the assertion also for enlarged control nets.

For completeness the time-modulus recurrence is explicit. Put

\[
a_j^z=\lambda S^2V_j^0/4,\quad a_j^h=sa_j^z,\quad
a_j^W=\lambda S^2\tau_jH_{j-1}/4\ (j\ge2),\quad
a_j^\delta=\lambda T_j^0/8,
\]

\[
a_L^k=\lambda H_L/4,\quad
a_j^k=a_{j+1}^W\tau_{j+1}+10a_{j+1}^\delta,\quad
A_t=2\mathcal K+(\lambda/4)SE_0,
\]

\[
v_t=A_tv_0+A_0(2\Gamma\tau_*+\lambda E_0/4)
                              +(\lambda/4)Sf_*B_0.
\]

Take $t_3=4\beta^2$, valid on the joining strip by Cauchy's formula, and
define

\[
z_1^t=Sv_t+D_c,\quad h_1^t=sz_1^t+t_2a_1^zz_1^0,
\]

\[
z_j^t=S(H_{j-1}v_t+a_{j-1}^hv_0)+a_j^Wh_{j-1}^0
                 +10h_{j-1}^t+D_c,\quad
h_j^t=sz_j^t+t_2a_j^zz_j^0,
\]

\[
k_L^t=v_t,\quad
k_j^t=a_{j+1}^Wd_{j+1}^0+10d_{j+1}^t
             +S\tau_{j+1}v_t+Sa_{j+1}^\delta v_0+D_c,
\]

\[
d_j^t=sk_j^t+t_2a_j^zk_j^0
 +(t_3M_0a_j^z+t_2a_j^k)z_j^0+t_2M_0z_j^t.
\]

The lower-Jacobian probes have coefficients

\[
z_1^J=0,\quad h_1^J=t_2a_1^zP_1,\quad
z_j^J=a_{j-1}^h+a_j^Wf_{j-1}+10h_{j-1}^J,\quad
h_j^J=sz_j^J+t_2a_j^zP_j.
\]

If $C_t$ is the maximum of $v_t,z_j^t,h_j^t,k_j^t,d_j^t,z_j^J,h_j^J$,
all terminal derivatives are bounded by
$C_t\sqrt{un}\,n^{\kappa_0}\ell_n^4$.
These formulas result by differentiating each displayed linearized product;
a raw coordinate derivative costs one $\sqrt n$, and each homogeneous
layer propagation remains $10s$. No derivative of the variational Jacobian
is needed in $v_t$, since its equation is differentiated only in time.
For a root block selector $E_i$, the quadratic matrix is
$\operatorname{sym}(E_i^\top L)$, so it has the same norm and moduli as
$L$. Centered interpolation costs at most $10u\|\Delta L\|$, including
the trace change, on $\|g\|\le2\sqrt{2u}$.

Here is a numerical bound for all these recurrences. The power bounds in
\Cref{cp:src-powers} and their geometric sums give

\begin{center}\small
\begin{tabular}{@{}p{0.57\linewidth}p{0.39\linewidth}@{}}
\toprule
Coefficients & Upper bounds\\
\midrule
$H_j,P_j,f_j$; $k_j,\tau_j$; $V_j^0$ & $X^3$; $X^4$; $X^7$\\
$A_*,D_*,\eta_*^{-1},\mathcal B,\Gamma,J_0$ & $X^5,X^6,X^{27},X^3,X^7,X$\\
$M_0,A_0,B_0,E_0$ & $X^{31},X^{33},X^{32},X^{38}$\\
$T_j^0,D_c$ & $X^{45},\lambda X^{45}$\\
parameter, forward, backward map coefficients & $X^{80}\chi,X^{85}\chi,X^{125}\chi$\\
control differences; reference speeds divided by $\lambda$ & $X^{90}\chi$; $X^{47}$\\
parameter, forward, backward time coefficients & $X^{125}\chi,X^{140}\chi,X^{190}\chi$\\
\bottomrule
\end{tabular}\end{center}

For example substitute the first two rows in $v_0$; its largest product
has exponent less than $80$ and at most one $T_0$ factor. The forward
recursion costs $P_j\le X^3$. In the backward recursion the extra $M_0$
is additive and its complete propagation costs
$(10s)^L\le X^2$, plus at most $L\le X$ summands. In the largest time
forcing $t_2M_0z_j^t$ the exponents sum to at most $1+31+140$; the
same propagation stays below $190$. These checks also cover the lower
probes. Thus a common coefficient $C_0=X^{300}\chi$ suffices.

The response maps require additional algebraic products, not a new flow.
For training driver $a$ and declared evaluated input $b$, put
$\bar R=R/S$ and $\bar Q=\phi'\odot\bar R$. Then

\[
\bar R_{ba}^{(1)}=\bar\delta_a^{(1)}v_a^\top v_b,\qquad
\bar R_{ba}^{(j)}=\bar\delta_a^{(j)}c_{ab}^{(j-1)}
                       +W^{(j)}\bar Q_{ba}^{(j-1)},\quad
c_{ab}^{(j-1)}=h_a^{(j-1)\top}h_b^{(j-1)}/n.
\]

Its linearized terms are $\bar\delta_{[1]}c+\bar\delta c_{[1]}$,
$SV_{\sqrt nW}\bar Q/\sqrt n$, $W\bar Q_{[1]}$ and the omitted
response port, where
$c_{[1]}=(h_{a,[1]}^\top h_b+h_a^\top h_{b,[1]})/n$ and
$\bar Q_{[1]}=\phi'\bar R_{[1]}+\phi''\bar Rz_{b,[1]}$.
Below a deletion retain the direct observable
$C_bC_a^\top\bar q_a$, $C_a=D_\Theta h_a$, and its forward images of
$C_a^\top y_i$. Its subtraction is

\[
(C_b^1C_a^{1\top}-C_b^0C_a^{0\top})y_i
=(C_b^1-C_b^0)C_a^{0\top}y_i
   +C_b^1(C_a^1-C_a^0)^\top y_i.
\]

The first term uses reference-probe coordinate localization; the second
uses the lower-adjoint bound in \hyperref[cp:src-local-lemma]{the local-map proof}.
Hence both retain the factor $d_0+n^{-1/2}$.
Differentiating the response products introduces
$\phi''\dot z_b\bar R$; use the coordinate response cap and RMS speed
of $z_b$, then convert to a scalar control once. With raw amplitude
coefficient at most $X^{80}$ and $U_{\rm fin}\le X^{26}$, these formulas
and geometric propagation give the common envelope

\[
X^{100}(1+X^{80}+4U_{\rm fin})^2(1+C_0)^2
 \le X^{1000}\chi^2=:\mathcal C.
\]

The exponent count is $100+160+600$, with the remaining $140$ paying
all numerical sums and bounded layer propagation. Thus the map norms,
control moduli and actual analytic derivatives are bounded by

\begin{equation}
\mathcal C\sqrt u n^{\kappa_0}\ell_n^5,\qquad
\mathcal C\sqrt u n^{\kappa_0}\ell_n^4,\qquad
\mathcal C\sqrt{un}n^{\kappa_0}\ell_n^8.
\label{cp:poly:eq:map-interface}
\end{equation}

There are at most $4u(N+1)^2$ real scalar controls, of amplitude
$M=\mathcal C\ell_n$ and speed $D=\mathcal C\sqrt n\ell_n^3$.
Every bound concerns frozen one-dimensional histories until the uniform
event is established.

We verify the geometry before forming that event. For the training
rectangle, the radius in \eqref{cp:poly:eq:training-radius} is
at most
$\min\{1/8,1/\lambda,\log2/(4\mathcal K),1/(32YSD_W)\}$
and at most $a/(64YSU_{\rm fin}\sqrt{\ell_n})$, by the source power
bounds. The residual, activity, operator and pole estimates in
\hyperref[cp:src-whole-sphere-verification]{the working-domain verification} therefore apply
without angular fields. The normalized proof has no $n^{-1}\le Y$ gate.

For the panel set $\nu=\lambda/4$ and

\[
h(t)=\frac1{2\mathcal K}\log(1+2\mathcal Kr_ne^{\nu t}),\quad
d_h=\nu/(2\mathcal K)\le1/8,\quad T_+=T+h(T)+r_n,
\]

\[
D_s=\{x+iy:-sh(0)\le x\le sT_+,\ |y|\le s h(x_+)\},
\]

\[
P_s(x+iy)=X_s(x)+i\operatorname{clip}_{[-s h(X_s(x)_+),s h(X_s(x)_+)]}y,
\quad X_s(x)=\operatorname{clip}_{[-sh(0),sT_+]}x.
\]

The retraction is $(1+d_h)$-Lipschitz and commutes with real anchors.
Its bounding rectangle has sides at most $H_0=128\ell_n/\lambda$;
canonical real-then-vertical contours have length at most $H_0/2$.
The derivative recurrence in \Cref{cp:src-time-derivative} gives
$J_*\le X^{34}$, and \eqref{cp:src-gradient-derivative} gives
$C_g\le X^{38}$; both are below $X^{100}$.
The complete propagator gate is the sum

\[
4\mathcal Kr_n+(2C_g/\sqrt{\mathcal K})B_nYr_n
 \le4\mathcal K/X^{300}+C_g/(8X^{300})<1/100.
\]

Here $Y=\lambda S/16$ and
$\lambda B_n/(1+\lambda B_n)\le1$. As in
\eqref{cp:src-vertical}, vertical residual activity is at most $Yr_n$;
freezing the real symmetric Gram generator makes its imaginary-time
evolution unitary, and the integrated generator difference is the second
term of this gate. Forward real segments contract. Thus ordered products
of base propagators cost less than two. Additional activity is at most
$8Yr_n$, normalized operator motion at most $8YSD_Wr_n$, and improved
vertical preactivation displacement at most
$2YSU_{\rm fin}\sqrt{\ell_n}r_n\le a_0/32$.
The negative corner costs at most $a_0/8$. All stopped margins are strict.

We detail the uniform endpoint event, because an endpoint clamped by a
full-network stop would not have frozen Gaussian coefficients.
In what follows $A=A_{N,\rho}=\mathcal C^4(N+d+u+1)^8$ and

\begin{equation}
A\ell_n^{16}\le n^{1/1000},\quad
u\le n^{1/8000},\quad \ell_n\le n^{1/16000}.
\label{cp:poly:eq:absorption-gates}
\end{equation}

To verify the first inequality at the source width put
$v=\log(A/\rho)\ge8000\log10$. At $\log n=40000v$,
$\log A+16\log(1+40000v)\le17v<40v$; the difference increases
thereafter. The other inequalities follow from the same gate.

For a cavity's own stop $\sigma_c$, condition only on its retained
initialization. Put $\xi=n^{-2}/(1+A\ell_n^8)$,
$J_s=\lceil16H_0/\xi\rceil$, $s_j=j/J_s$, and take a rectangular
product grid $\mathscr G$ with side spacing at most $\xi/8$ on the
bounding rectangle. Test all endpoints

\begin{equation}
z_{j,q,c}=P_{\min(s_j,\sigma_c)}q,
\qquad q\in\mathscr G,\quad0\le j\le J_s.
\label{cp:poly:eq:prefix-grid}
\end{equation}

Each endpoint is retained-measurable. Its contour is explicitly
$\zeta(t)=\operatorname{sgn}(x)\min(t,|x|)
+i\operatorname{sgn}(y)\min((t-|x|)_+,|y|)$ on $[0,H_0]$, followed
by freezing. It lies in the cavity's own domain. On this contour solve
$V'=\zeta'(\mathsf AV+\sum_a\mathsf B_a^ee_a+\mathsf B_a^q\bar q_a)$,
$V(0)=0$, where the coefficients are the differentiated port forces above.
Bounded measurable histories suffice; this equation is a linear map of
the omitted roots. No clamped extension is claimed holomorphic.

For control accuracy $\tau=n^{-1/8}$ sample $[0,H_0]$ in
$\lceil16DH_0/\tau\rceil$ equal intervals, round each real component
to mesh $\tau/8$ and interpolate. Every admissible history is within
$\tau$ and the logarithmic cardinality is at most

\[
4u(N+1)^2(2+16DH_0/\tau)\log(1+32M/\tau)
 \le A\ell_n^8n^{5/8}\le n^{0.627}.
\]

The coefficient is paid because
$16384u(N+1)^2\mathcal C\lambda^{-1}\le A$.
The net uses a loose amplitude bound $2\mathcal C\ell_n$.
Affine dependence on controls and \eqref{cp:poly:eq:map-interface}
therefore bound every map on the entire net by
$4\mathcal C^2\sqrt u n^{\kappa_0}\ell_n^5\le n^{1/800}<n^{1/200}$.

For $g\sim N(0,I_r/n)$, $r\le2un$, a real symmetric matrix obeys

\[
\Pr\{|g^\top Rg-\operatorname{tr}R/n|>t\}
 \le2\exp\left[-\frac18\min\left\{
 \frac{n^2t^2}{\|R\|_{\rm HS}^2},\frac{nt}{\|R\|}\right\}\right].
\]

Diagonalization and
$-\log(1-v)-v\le v^2/[2(1-|v|)]$, followed by exponential Markov
with parameter at most $n/(4\|R\|)$, prove this bound.
Since $\|R\|_{\rm HS}^2\le2un\,n^{1/100}$, real/imaginary splitting
at threshold $n^{-1/10}/8$ gives at most $4e^{-n^{0.78}}$.
The linear-coordinate tail is smaller. The Euclidean image mean is at
most $\sqrt{2u}n^{1/200}<n^{1/100}/4$ and Gaussian norm concentration
pays its remaining threshold $n^{1/100}/2$. The root-norm failure is
at most $e^{-un}$, and control interpolation costs at most
$10\mathcal C u^{3/2}n^{\kappa_0-1/8}\ell_n^4<n^{-1/10}/8$.

Now select any common prefix $D_\sigma$, even one depending on all roots.
For $z\in D_\sigma$ choose a nearest grid point $q$ and the largest
$s_j\le\sigma$. Clipping gives
$\sup_z|P_sz-P_tz|\le3H_0|s-t|$; hence

\[
|z-P_{s_j}q|\le(3/16+9\sqrt2/64)\xi.
\]

Its canonical contour lies inside $D_\sigma$, so actual controls may now
be substituted. With those actual analytic functions fixed, the linear
ODE has a unique holomorphic solution on this simply connected vertically
fibered domain. Local power-series uniqueness and continuation along compact
paths prove this assertion; contraction to the real interval gives simple
connectivity. The image under $P_\sigma$ of the straight joining segment
has length below $\xi/2$. Integrating the actual derivative bound in
\eqref{cp:poly:eq:map-interface} along it costs at most
$10u\mathcal C\sqrt u n^{\kappa_0+1/2}\ell_n^8\xi<n^{-1/10}/8$.
Boundary levels follow by continuity; level zero is a tested endpoint.
The endpoint family has at most
$(2+16H_0/\xi)^3\le n^9$ members. Together with at most $An^3$ maps
per endpoint and $uLn^u$ deletion sets, the total logarithmic union cost
is below $n^{0.65}$, so its failure is below
$e^{-n^{0.7}}<\rho/16$.

For the training rectangle an even simpler grid works: round each real
coordinate toward zero to a multiple of $\xi$. The resulting grid point
and its straight segment to the target lie in every common rectangle
prefix. At each fixed grid point first clamp the canonical contour only
by the cavity's own rectangle. The same conditional tests and interpolation
then apply, with at most $n^6$ endpoints. Thus the preceding larger
probability allowance covers both geometries.

The exact nonlinear identities \eqref{cp:src-exact-remainders},
\eqref{cp:src-remainder-bounds} and
\eqref{cp:src-force-remainder} now close the comparison quantitatively.
Take $N_0=n^{1/100}$, $d_0=n^{-1/10}$ and $u_0=n^{-1/25}$ for the
linear Euclidean, linear coordinate and nonlinear remainder stops.
Their variables are $R=d_0N_0+d_0u_0+u_0^2$ and
$P=(N_0+u_0)^2/\sqrt n$.
The force bound has coefficient $X^{240}(1+u)(1+M)$, and learned
ports and the top offset have coefficient $X^{260}(1+u)^2(1+M)^3$.
Inserting $M\le X^{80}\ell_n$, the map coefficients above, activity
$2\int\rho/S\le1$ and length $64\ell_n/\lambda$ gives

\[
\|U\|\le A\ell_n^8n^{\kappa_0}
 [n^{-0.08}+u_0^2+n^{-0.48}+n^{-1/2}]<u_0/2.
\]

The largest resulting power is $n^{-0.07875}$, below $u_0/2$ at the
source gate. The joining-strip inequality
$d_0+X^{10}u_0+X^{30}(R+P)<a/32$ holds at the same gate.
This is a first-exit argument. It gives retained coordinate differences
below $n^{-1/30}$ and vector differences below $2n^{1/100}$, including
the augmented responses. At the first layer the incoming-row increment
is bounded by $sS^2M$, with no $n^{-1/2}$; it enters the deleted scalar
controls, not a nonexistent retained reverse port.

The four nonzero scalar traces are exactly
\eqref{cp:src-traces}. The normalized Schatten calculation and scalar
feedback solution in \hyperref[cp:src-moments]{the trace calculation} apply
before adaptive control substitution. Their first three traces keep the
original constants; the adaptive-residual trace is at most
$A\ell_n^8n^{-1+\kappa_0}$, and finite remainders are below $n^{-1/40}$.
Thus the original $C_{\rm abs}$ works. A finite Gaussian grid improves
training maxima to $K_{\rm src}\sqrt{\ell_n}$ and
$SK_{\rm src}\sqrt{\ell_n}$. For the response the four deterministic
coefficients, after division by $S$, are

\[
sK_{\rm src}H_{j-1}^2,\quad sK_{\rm src}f_{j-1}^2,\quad
sK_{\rm src}ST_Q,\quad sK_{\rm src}S^2H_{j-1}q_{j-1}^*.
\]

The centered row coefficient is $64q_{j-1}^*$, giving precisely
$U_j(64)$. There are at most $mNLn$ response tests times the polynomial
endpoint count; Gaussian multiplier $64$ pays them at the source gate.
Both maximum failures together are below $\rho/16$.
The exact velocity $\dot z_b=-(2/m)\sum_a r_aR_{ba}$ then improves
poles and passive maxima. The real displacement is at most
$2S^2U_{\rm fin}\sqrt{\ell_n}$; no passive budget is required.

It remains to remove budgets at the actual confidence order $u$, rather
than take a limit at each fixed order. Complete each common-cavity
reference by its own stopped retraction before integrating a root.
For its normalized complex-minus-real Gaussian coefficients the panel
derivative estimate gives

\[
D\le X^{100}\lambda r_nB_n,\quad K\le X^{100}\lambda B_n,
\quad H\le128\ell_n/\lambda.
\]

Use the larger radius
$\bar D=X^{-200}\lambda B_n/[(1+\lambda B_n)\sqrt{b_n}]$.
The rectangle lemma \Cref{cp:src-Gaussian} bounds the logarithmic
moment by $a(M_G+2\bar D)+2a^2\bar D^2$, where
$M_G=256\bar D\sqrt{\log(e+HK/\bar D)}$.
Here

\[
\log(e+HK/\bar D)\le300\log X+16b_n+\log(1/\lambda).
\]

For $b\ge1,x\ge0$,
$\min(1,be^{-x})\sqrt x\le\sqrt{\log b}+1$, by splitting at
$x=\log b$ and writing the excess as $s$ with $e^{-s}\sqrt s\le1$.
Apply this with $b=B_n$, $x=\log(1/\lambda)$ and use
$B_n\le1+b_n$. Then
$M_G+2\bar D\le X^{-190}$, so every such correction has moment
at most $e^{X^{-180}}$ for $0\le a\le8$, including its sample RMS.

For the smaller training rectangle use instead
$D\le X^{40}\lambda B_nr$, $K\le X^{40}\lambda B_n$,
$H\le64\ell_n/\lambda$. Substitution of
\eqref{cp:poly:eq:training-radius} bounds the mean plus $2D$
by $X^{-50}$. To check uniformity, split the entropy into
$\log128+\log u+100\log X+\log\chi+(3/2)\log\ell_n$ and use
$u^{-1}\sqrt{\log u}\le1$, $\lambda\sqrt{\log\chi}\le2$, and
boundedness of
$(1+\log(e+\ell_n))(1+\sqrt{\log\ell_n})/\sqrt{\ell_n}$.
The numerical bound for the last expression follows from
$\log(e+x)\le2+\log x$ and maximizing each
$(\log x)^a/\sqrt x$, $a\in\{0,1/2,1,3/2\}$.
Five of the sixty available powers of $X$ absorb the constants.
Thus fixed-exponent corrections are small on both domains.

For each distinct omitted root $i$ form its \emph{complete} common-cavity
reference weight $A_i$: exponentiate $\eta_*C_{\rm abs}$ times one
plus its four real scalar/sample-RMS references and four complex-minus-real
corrections. The real moment \eqref{cp:src-main-moment} and
Cauchy--Schwarz give $\mathbb E(A_i^2\mid\text{cavity})<3$;
Hölder within the four correction terms uses exponent
$16\eta_*C_{\rm abs}<8$. Different $A_i$ are conditionally independent
because they use distinct root pairs of the same cavity.

Only the singleton-to-common-cavity differences require order $u$.
Extend both paths by their separate own stops and radially project their
whole difference to the ball of radius $4n^{1/100}$. Both omit the paired
root. Their Gaussian radius is $v_n=4n^{-49/100}$ and the same rectangle
lemma gives mean at most

\[
\mu_n=256v_n\sqrt{\log(e+X^{400}\chi\ell_n^4/v_n)}+2v_n.
\]

The source gate gives $u\mu_n+u^2v_n^2<10^{-3}$.
For $B_i^{\rm wt}=e^{\eta_*(Z_{a,i}+K_{a,i}/S)}$, scalar absorption,
Cauchy--Schwarz between main references and corrections, and Hölder over
at most four corrections per root give, for distinct $I$, $|I|=k\le u$,

\[
\mathbb E\left[1_E\prod_{i\in I}B_i^{\rm wt}\right]
 \le3^{k/2}\exp\{kn^{-1/40}+4k\mu_n+64k^2v_n^2\}\le16^k.
\]

The event indicator is dropped only after pathwise domination by complete
references. No independence of sample processes or correction paths is used.
From the defining constants,
$2\eta_*K_{\rm src}\le33/4096<1/100$, so repeated weights cost at most
$W_n=e^{\sqrt{\ell_n}/100}$. Partition an ordered $u$-tuple by its
distinct indices. A partition with $j$ repetitions is encoded by $j$
edges from block minima, giving at most $\binom u2^j/j!$ possibilities.
The preceding distinct-product bound therefore yields

\[
\mathbb E\left[1_E(n^{-1}\sum_iB_i^{\rm wt})^u\right]
 \le16^u\exp\{\tbinom u2W_n/(16n)\}\le e16^u,
\]

where $n\ge4u^3$ suffices and is implied by the source gate.
Markov and the sample/layer union make the budget-hit probability at most
$emL(64e^2)^{-u}\le\rho/4$.
The transferred cavity budgets obey
$\mathcal H_a^{\rm cavity}\le e^{2\eta_*n^{-1/30}}\mathcal H_a+u/n
<2\mathcal B$. All other stops have the strict margins already checked.
For either geometry the four failure allowances are respectively
$\rho/4$, $\rho/16$, $\rho/16$ and $\rho/4$.
For the simultaneous rectangle and panel event, initialization is shared;
sum each of the other three allowances twice. The total is strictly below
$\rho$. Holomorphic local existence and uniqueness then continue the
solutions through the compact domain boundary. This completes the finite
probability calculation.

The all-time real extension needs no new probability. Fitting gives
$\|\theta(t)-\theta(T)\|\le2Y\lambda^{-1/2}e^{-16\ell_n}$ in
Euclidean mobility coordinates. The backward-only subtraction in
\hyperref[cp:src-finite-working]{the finite-query proof} costs at most $X^{20}n$ and gives
the stated carrier cap. Subtract the response recursion with the state at
$T$ fixed. Changed gates multiply only reference coordinate caps, while
changed mixers and pairings use operator and RMS bounds. A coefficient
$X^{100}\chi^2(1+\sqrt{\ell_n})$ bounds this RMS Lipschitz subtraction;
one conversion to coordinates leaves a tail at most
$YX^{103}\chi^3(1+\sqrt{\ell_n})n^{-31/2}$.
The source gate makes it at most $Y\chi\sqrt{\ell_n}$, proving
\eqref{cp:poly:eq:real-response}.

Finally the panel disk radius is $\varrho(t)=h(t)/[2(1+d_h)]$.
As in \hyperref[cp:src-panel-working]{the panel-domain proof}, it lies inside the domain and
the partition $t_{b+1}=\min(T,t_b+\varrho(t_b)/2)$ satisfies

\[
J\le1+40/(\lambda r_n)
 +80\mathcal K\lambda^{-1}\log(1+4\ell_n/\log2).
\]

The reciprocal minimum in \eqref{cp:poly:eq:panel-radius} is at most

\[
\frac{1024U_{\rm fin}}{a_0}(Y/\lambda)^2\sqrt{\ell_n}
 +X^{300}\lambda^{-1}\sqrt{b_n}+X^{300}B_n\sqrt{b_n}.
\]

For $x=\log\ell_n\ge0$, $b_n\le2+x$ and
$(2+x)^{3/2}e^{-x/2}\le4$. Hence $b_n^{3/2}\le4\sqrt{\ell_n}$.
Use $S^2=256(Y/\lambda)^2$, $a_0^{-1}\le\beta$ and the coefficient
bounds to obtain \eqref{cp:poly:eq:panel-count}.
This proves the source lemma.

\end{proof}

\subsection{A finite lower bound for actual dense variability}
\label{cp:poly:dense-lower-section}

For $0<\tau<1$, define

\[
a_\tau=\frac{\tau\gamma^2}
 {2^{20}m u_{\tau/8}X^{108}(1+\lambda)},
\]

\begin{equation}
N_{\rm low}(\tau)=\left\lceil\max\left\{\begin{gathered}
 N_{\rm src}(m,\tau/8),\quad
 2^{44}m^6X^{96}\gamma^{-12}\tau^{-2}\log(64/\tau),\\
 [2^{30}u_{\tau/8}X^{114}m\chi(1+\gamma^{-1})^2/\tau]^2
 \end{gathered}\right\}\right\rceil.
\label{cp:poly:eq:lower-width}
\end{equation}

\begin{proposition}[Quantitative trajectory separation]
\label{cp:poly:dense-lower}

For each $n\ge N_{\rm low}(\tau)$, two independent dense runs satisfy

\begin{equation}
\Pr\left\{\|f_n-\widetilde f_n\|_{\mathcal X}
       \ge a_\tau\frac{Y}{\sqrt n\ell_n^{5/2}}\right\}\ge1-\tau.
\label{cp:poly:eq:dense-lower}
\end{equation}

The domain may be any promised query set containing the training inputs.
The lower coefficient is weaker than that in \eqref{cp:eq-dense-lower}.

\end{proposition}

\begin{proof}

We first prove a shift-uniform finite scalar inequality. For iid centered
real variables $V_i$ with variance $\sigma^2>0$ and third absolute moment
at most $M_3$, put $Z=n^{-1/2}\sum_iV_i$. For every $h>0$ and real $b$,

\begin{equation}
\Pr\{|Z-b|\le h\}\le3h/\sigma
 +4\exp[-9h^2\sigma^4n/(128M_3^2)].
\label{cp:poly:eq:finite-smallball}
\end{equation}

For an independent copy $V'$, Taylor expansion gives
$|\mathbb Ee^{itV}|^2=\mathbb E\cos(t(V-V'))
\le1-\sigma^2t^2+(4/3)M_3|t|^3$, because
$\mathbb E|V-V'|^3\le8M_3$. Therefore
$|\mathbb Ee^{itV}|\le e^{-\sigma^2t^2/4}$ for
$|t|\le3\sigma^2/(8M_3)$.
The Gaussian majorant
$1_{|z-b|\le h}\le\sqrt e\,e^{-(z-b)^2/(2h^2)}$ and its Fourier
integral bound the probability by

\[
\frac{\sqrt e\,h}{\sqrt{2\pi}}
 \int_{\mathbb R}e^{-h^2t^2/2}|\mathbb Ee^{itZ}|\,dt.
\]

On $|t|\le3\sigma^2\sqrt n/(8M_3)$ use the preceding Gaussian
bound; its full-line integral is below $3h/\sigma$.
On the complement use modulus one and
$\Pr(|G|>s)\le2e^{-s^2/2}$. This proves
\eqref{cp:poly:eq:finite-smallball}. Fourier integrability justifies the
exchange of expectation and integration. The argument works conditionally
on any sigma-field supplying conditional iid variables and moment bounds;
the shift may be measurable in that sigma-field.

Condition on the first run through layer $L-1$ and the entire independent
second run. Its final feature rows $U_i\in\mathbb R^m$ are conditionally
iid. The initialized covariance calculation above, now choosing

\[
e_1=\min\{(8T_1)^{-1},
 \gamma/[4m\max(C_0^{\rm cov},s^2C_0^{\rm cov}+D_0^{\rm cov}T_1)]\},
\]

gives, except with probability at most $\tau/8$,

\[
\bar Q:=\mathbb E(U_iU_i^\top\mid\mathcal F)\succeq(\gamma/2)I,
\qquad \max_a\|U_{i,a}\|_{L^6(\cdot\mid\mathcal F)}\le X^4.
\]

Indeed the conditional last-layer mean entry error is at most
$s^2E_{L-1}+D_0^{\rm cov}\Delta_{L-1}\le\gamma/(4m)$.
Its required width is
$32s^2R_0^2e_1^{-2}\log[256L(m+1)^2/\tau]$, which is bounded by
the source gate in \eqref{cp:poly:eq:lower-width}. Linear growth
and Gaussian sixth moments prove the $L^6$ bound.

The squared-area calculation in \hyperref[cp:innovation]{the innovation proof} does not
require equal diagonal moments. For any random $U$ with
$\bar Q\succeq gI$, put $c=\bar Qy$, $T_q=\operatorname{tr}\bar Q$,
$M_4=\mathbb E\|U\|^4$, and
$V=\mathbb E\|U(U^\top y)-c\|^2$. It proves

\[
V\ge\frac{[y^\top\bar Q^2(T_qI-\bar Q)y]^2}
 {4M_4\|\bar Qy\|^2}
 \ge\frac{(m-1)g^3(T_q-g)}{4M_4}\|y\|^2.
\]

For clarity the last step multiplies the two spectral lower bounds
$(m-1)g\|\bar Qy\|^2$ and $g^2(T_q-g)\|y\|^2$ for the numerator
before squaring. Every eigenvalue lies in $[g,T_q-(m-1)g]$, which
verifies both bounds. With $g=\gamma/2$, $T_q\ge mg$ and
$M_4\le m^2X^{16}$, at least one coordinate satisfies

\[
\operatorname{Var}[U_a(U^\top y)\mid\mathcal F]
 \ge\sigma_0^2:=Y^2\gamma^4/(256X^{16}).
\]

Choose the smallest maximizing index, which is $\mathcal F$-measurable.
The centered product has third absolute moment at most
$8m^3Y^3X^{24}$: use $\sum_b|y_b|\le mY$, the conditional $L^6$
bound, and $\mathbb E|V-\mathbb EV|^3\le8\mathbb E|V|^3$.
Thus $M_3^2/\sigma_0^6\le2^{30}m^6X^{96}\gamma^{-12}$.
Zero readout gives the exact identity

\[
\dot f_n(0,x_a)=\frac2{mn}\sum_iU_{i,a}(U_i^\top y).
\]

The second run and the conditional means are only a measurable shift.
Apply \eqref{cp:poly:eq:finite-smallball} with
$h=\tau\sigma_0/24$; its scalar failure is at most $\tau/4$ under
the second gate of \eqref{cp:poly:eq:lower-width}. Hence, allowing
the covariance failure, with probability at least $1-\tau/2$,

\[
\max_a|\dot f_n(0,x_a)-\dot{\widetilde f}_n(0,x_a)|
 \ge\frac{\tau Y\gamma^2}{192mX^8\sqrt n}.
\]

Use the training rectangle with failure $\tau/8$ for each run. The
parameter-two ellipse on $[0,r]$ lies inside it, and the prediction
difference is bounded by $32X^6Y/\lambda$. Taking
$K=\lceil2\ell_n/\log2\rceil\le4\ell_n$ in
\Cref{cp:derivative-transfer} yields

\[
\|f_n-\widetilde f_n\|_{\mathcal X}
 \ge\frac{\tau Y\gamma^2}
 {6144m u_{\tau/8}X^{108}(1+\lambda)\sqrt n\ell_n^{5/2}}
       -\frac{768X^6Y}{\lambda n^2}.
\]

The last term is at most half the first whenever
$n^{3/2}\ge2^{24}u_{\tau/8}X^{114}
m(1+\lambda)\ell_n^{5/2}/(\lambda\tau\gamma^2)$.
Since $\ell_n^{5/2}\le64n^{1/4}$, the third gate in
\eqref{cp:poly:eq:lower-width} suffices. The probability union
costs less than $\tau$ and the stated coefficient is smaller than the
remaining leading coefficient. Both predictors start at zero, so the
witness time is strictly positive. At training inputs their fitted
limits agree with the same label; no endpoint separation follows.

\end{proof}

\subsection{The third-order Legendre history estimate}
\label{cp:poly:legendre-section}

The Legendre model is exactly \eqref{cp:eq-legendre-reconstruction} and
\eqref{cp:eq-legendre-flow}, with its constant forward and zero backward
prefixes. Let

\begin{equation}
q=\left\lceil\lambda^{-3}n^{1/4}\ell_n^2e^{\sqrt{\ell_n}/2}\right\rceil,
\qquad P=X^{500}\chi^{10}.
\label{cp:poly:eq:legendre-order}
\end{equation}

For fixed positive $a,b$, define the numerical logarithm coefficient

\[
C(a,b)=\max\{1,e^b(a/(eb))^a\}.
\]

Maximizing $x^ae^{-bx}$ over $x\ge1$ proves
$\ell_n^a\le C(a,b)n^b$. Define the structural gate

\begin{equation}
N_{\rm Leg}=\left\lceil\max\left\{
(2e^{100}\lambda^{-3})^8, (4P^3)^4,
[8P^4C(2,1/64)]^8, d^4
\right\}\right\rceil.
\label{cp:poly:eq:legendre-gate}
\end{equation}

\begin{proposition}[Improved error without additional state]
\label{cp:poly:legendre}

On the finite training-source event, for $n\ge N_{\rm Leg}$,

\begin{equation}
\|f_{\rm Leg}-f_n\|_{\mathcal X}
 \le4YP^5\ell_n^2n^{-47/64}.
\label{cp:poly:eq:legendre-error}
\end{equation}

The exact moving and fixed inventories remain
$n(d+1)+1+2(L-1)mnq$ and $(L-1)n^2$, respectively.

\end{proposition}

\begin{proof}

Use Euclidean mobility coordinates
$\theta=(W^{(1)}/\sqrt n,W^{(2)},\ldots,W^{(L)},w/\sqrt n)$.
For a finite physical horizon $T_0$ put

\[
d_{T_0}=\int_0^{T_0}\sum_{j=2}^L\|\mathcal E_j(t)\|_Fdt,\qquad
D_{T_0}=\sup_{t\le T_0}\sum_{j=1}^{L+1}
                  \|\widehat\theta_j(t)-\theta_j(t)\|,
\qquad A=\tau(T_0).
\]

Here $\mathcal E_j$ is the exact physical defect in
\eqref{cp:eq-legendre-defect}. The order-independent fitting lemma gives

\[
\widehat\rho(t)\le Ye^{-\lambda t/4},\qquad
\int_t^{T_0}\widehat\rho(s)ds\le4\widehat\rho(t)/\lambda,
\qquad A\le1+4Y/\lambda\le2.
\]

Work first on the bootstrap $D_{T_0}\le Y/\sqrt n$.
Subtract backward recursions with the dense trajectory as reference, so
changed gates multiply only its known coordinate carrier. The all-time
response bound \eqref{cp:poly:eq:real-response} and the exact
recursion for $R$ then give

\[
\max_{a,j}\|\widehat k_a^{(j)}-k_a^{(j)}\|_2/\sqrt n
 \le X^{40}(1+\sqrt{\ell_n})D_{T_0},
\]

\[
\max_{a,b,j}\|\widehat R_{ba}^{(j)}-R_{ba}^{(j)}\|_2/\sqrt n
 \le X^{130}\chi^4(1+\sqrt{\ell_n})D_{T_0}.
\]

The nontrivial products can be checked directly: forward differences cost
$X^5D$, feature-pair differences $X^9D$, and the changed response gate
costs at most $YX^{116}\chi^3(1+\sqrt{\ell_n})D$ using the reference
response maximum. Homogeneous propagation costs $(9s)^L\le X^2$,
with at most $L\le X$ terms. Thus the displayed powers include all terms.
One RMS-to-coordinate conversion and the bootstrap yield

\[
\max|\widehat k|\le YX^{50}\chi(1+\sqrt{\ell_n}),\qquad
\max|\widehat R|\le YX^{135}\chi^4(1+\sqrt{\ell_n}).
\]

Let $\widehat v_a=\widehat r_a/\widehat\rho$ and let $\mathcal F$
be the ordinary dense vector field evaluated at the reconstructed closure
state. At finite times $\widehat\rho>0$ by uniqueness in the fitting
proof; stationary continuation covers residual zeros. The normalized
ordinary feature velocity is

\[
a_b^{(j)}(t)=D_\theta h_b^{(j)}(\widehat\theta(t))
 \frac{\mathcal F(\widehat\theta(t))}{\widehat\rho(t)}
=-\frac2m\phi'(\widehat z_b^{(j)})\odot
                          \sum_a\widehat v_a\widehat R_{ba}^{(j)}.
\]

The residual calculation in the Legendre fitting proof gives
$\|\widehat v\|_2/\sqrt m=1$ and
$\|\dot{\widehat v}\|_2/\sqrt m\le X^{11}$.
Differentiating the \emph{actual} closure trajectory, including its defect,
gives the following bounds:

\begin{equation}
\begin{aligned}
\max_{b,j}\|\dot{\widehat z}_b^{(j)}\|_2/\sqrt n
 &\le YX^{30}\chi^3\widehat\rho,\\
\max_{a,j}\|\dot{\widehat\delta}_a^{(j)}\|_2/\sqrt n
 &\le X^{110}\chi^4(1+\sqrt{\ell_n})\widehat\rho,\\
\max_{a,b,j}\|\dot{\widehat R}_{ba}^{(j)}\|_2/\sqrt n
 &\le X^{220}\chi^8(1+\sqrt{\ell_n})\widehat\rho,\\
\max_{b,j}(\|a_b^{(j)}\|_2+\|\dot a_b^{(j)}\|_2)/\sqrt n
 &\le YX^{230}\chi^8(1+\sqrt{\ell_n}).
\end{aligned}
\label{cp:poly:eq:legendre-derivatives}
\end{equation}

To verify these estimates, \eqref{cp:eq-legendre-pointwise-defect}
bounds the full hidden block speed by
$YX^{20}\chi^3\widehat\rho$; forward differentiation costs $X^5$.
In $\dot\delta=\phi''\dot z\,k+\phi'\dot k$, use the coordinate
maximum of $k$ and RMS of $\dot z$. At the top
$\|\dot w\|/\sqrt n\le X^4\widehat\rho$; below it,
$\dot k=\dot W^\top\delta+W^\top\dot\delta$.
The forcing coefficient is at most
$X^{90}\chi^4(1+\sqrt{\ell_n})$, covered by the second line after
bounded layer propagation. Differentiating $R=\delta c+WQ$ adds
$\dot\delta c$, $\delta\dot c$, $\dot WQ$,
$W\phi''\dot zR$ and $W\phi'\dot R$.
The curvature product uses the coordinate maximum of $R$ and RMS of
$\dot z$; all remaining terms use pairing derivatives and operator bounds.
The largest forcing is below
$X^{180}\chi^8(1+\sqrt{\ell_n})\widehat\rho$, giving the third line.
Finally differentiate $a_b^{(j)}$: the three terms use $\dot z$ times
the response maximum, $\dot{\widehat v}$ times response RMS, and
$\dot R$ times $\|\widehat v\|_2/\sqrt m=1$.
Since $\widehat\rho\le Y$, they give the fourth line. In particular no
coordinate bound on $\dot z$ or $\sqrt n\|\mathcal E\|$ is introduced.

Subtract the accumulated physical feature defect from its history:

\[
g_b^{(j)}(\tau(t))=\widehat h_b^{(j)}(t)
 -\int_0^tD_\theta h_b^{(j)}(\widehat\theta(s))\mathcal E(s)\,ds,
\]

and use its constant initial value on $[0,1]$. Aggregate history norms
use $(mn)^{-1/2}$ over samples and neurons. The forward derivative bound
gives $\|h-g\|_{L^2}/\sqrt{mn}\le X^{10}d_{T_0}$.
Exactly $\partial_\tau g=a$ and
$\partial_\tau^2g=\dot a/\widehat\rho$ on $(1,A)$.
The first derivative vanishes at the prefix join because the initial
readout vanishes. Thus no distributional jump occurs there.

For $\mathscr D_Ag=-\partial_\xi[\xi(A-\xi)\partial_\xi g]$,
the inequality $(A-\tau(t))/\widehat\rho(t)\le4/\lambda$ cancels
the apparent terminal singularity. By
\eqref{cp:poly:eq:legendre-derivatives},

\[
\|\mathscr D_Ag\|_{L^2}/\sqrt{mn}
 \le YX^{240}\chi^9(1+\sqrt{\ell_n}).
\]

The shifted Legendre modes have eigenvalues $j(j+1)$.
Two integrations by parts have zero endpoint flux and no prefix jump;
Bessel's inequality then gives
$\|(I-\Pi_q^A)g\|_{L^2}\le[q(q+1)]^{-1}\|\mathscr D_Ag\|_{L^2}$.
Projection contraction on $h-g$ yields the forward tail

\[
H_{T_0}:=\max_j\|(I-\Pi_q^A)h^{(j)}\|_{L^2}/\sqrt{mn}
 \le YP\ell_n/q^2+Pd_{T_0}.
\]

The backward-history freezing calculation following
\eqref{cp:eq-legendre-ledger} is unchanged. Its coefficients satisfy
$b_0\le X^{42}(Y/\lambda)(1+(Y/\lambda)\sqrt{\ell_n})$,
$b_1\le X^{16}(Y/\lambda)\sqrt\chi$ and
$B_d\sqrt{4Y/\lambda}\le P\ell_n$.
Therefore, for the history
$b_a=(\widehat r_a/\widehat\rho)\widehat\delta_a$,

\[
B_{T_0}:=\max_j\|(I-\Pi_q^A)b^{(j)}\|_{L^2}/\sqrt{mn}
 \le YP\ell_n/q+P\ell_nD_{T_0}
 \le2YP\ell_n/q\qquad(q\le\sqrt n).
\]

For each fixed full history the growing-projection identity
\eqref{cp:eq-projection-growth} integrates its squared endpoint error
to its final squared tail norm, since initial errors vanish.
Applying Cauchy--Schwarz in samples and activity time to the exact defect
thus gives $d_{T_0}\le2(L-1)H_{T_0}B_{T_0}$.
The signed physical comparison \eqref{cp:eq-legendre-signed} and its
coefficient ledger give

\[
D_{T_0}\le\mathcal A_nd_{T_0},\quad
\mathcal A_n\le eX\exp[X^{36}(Y/\lambda)^2\sqrt{\ell_n}]
 \le X^2n^{1/64}\le Pn^{1/64}.
\]

For the last inequality use $X^{36}(Y/\lambda)^2\le X^{-24}$ and
$c\sqrt{\ell_n}\le\ell_n/64+16c^2$ with $c\le X^{-24}$.
Absorb the term linear in $d_{T_0}$ when $q\ge4P^3\ell_n$, using
$Y\le1$ and $2(L-1)\le P$. This yields

\begin{equation}
d_{T_0}\le4YP^3\ell_n^2/q^3,\qquad
D_{T_0}\le4YP^4\ell_n^2n^{1/64}/q^3.
\label{cp:poly:eq:third-order-defect}
\end{equation}

The gate \eqref{cp:poly:eq:legendre-gate} ensures
$q\le\sqrt n$, $q\ge4P^3\ell_n$, and
$8P^4\ell_n^2\le n^{15/64}$.
For the first use
$\ell_n^2e^{\sqrt{\ell_n}/2}\le e^{100}n^{1/8}$ and absorb the ceiling
by two. The second follows from $n^{1/4}\ge4P^3$.
For the third use $\ell_n^2\le C(2,1/64)n^{1/64}$.
Consequently the bootstrap improves to $D_{T_0}\le Y/(2\sqrt n)$.
Continuity excludes a first exit. All constants are independent of $T_0$,
so the estimate extends to all times and to the fitted limits. Real
whole-sphere prediction subtraction costs at most $X^7D_{T_0}$,
proving \eqref{cp:poly:eq:legendre-error}.
The exact inventory is \eqref{cp:eq-legendre-count}; $n\ge d^4$
absorbs the first-matrix term in the original headline count.

\end{proof}

\subsection{Taylor accuracy, storage and final width gates}
\label{cp:poly:assembly}

Use the same exact-metric selection and selected dynamics as
\Cref{cp:selected}. Its proof needs the real fitting, RMS, carrier,
and finite source-approximation bounds; the polynomial source lemma
supplies each of them on the capped panel domain.
The explicit comparison ledger permits $C_{\rm sel}=X^{60}$:
it gives $F\le X^6$, $K_1\le X^{18}$, $H_C\le2X^3$,
$F_c\le X^{15}$ and
$\mathcal B_n/(Y/\lambda)\le X^{22}(1+\sqrt{\ell_n})$.
Using $e^v-1\le ve^v$, its finite-time and tail coefficients are below
$C_{\rm sel}(Y/\lambda)(1+\lambda^{-1/2})(1+\sqrt{\ell_n})
e^{\sqrt{\ell_n}}$ and
$C_{\rm sel}(Y/\lambda)(1+\lambda^{-1/2})$, respectively.
Choose exactly the tolerance

\[
\eta=\min\left\{1,Y,16Y/\lambda,
\frac{\lambda}
 {2C_{\rm sel}n(1+\lambda^{-1/2})(1+\sqrt{\ell_n})e^{\sqrt{\ell_n}}}
\right\}.
\]

The finite-time error is at most $Y/(2n)$. The fitted-tail contribution
is at most $Y/(2n)$ under the label-independent gate

\begin{equation}
n^7\ge2C_{\rm sel}\lambda^{-1}(1+\lambda^{-1/2}).
\label{cp:poly:eq:taylor-tail-gate}
\end{equation}

Indeed $ne^{-8\ell_n}\le n^{-7}$.
The disk cover strictly contains the real interval in its open union, so
some positive closed rectangle lies inside it by compactness. Thus the
initialization-only compiler \Cref{cp:jets} applies without a new
setup-cost claim.

Let $k$ be the dimension of the declared input span, $M_n=X^6\sqrt n$,
and

\[
K=\max\{0,\lceil\log_2(8M_n/\eta)\rceil\},\qquad
R=2m+k+1+2(2m+p)J(K+1).
\]

On each half-radius panel the Taylor tail is at most $M_n2^{-K}\le\eta/8$.
Approximate each base coefficient in Euclidean norm to $\eta/128$ and
apply its initialized matrix exactly. Geometric summation gives base
coefficient error at most $\eta/64$ and image error at most $\eta/8$,
since initialized operators have norm eight. This is the same exact-pairing
source contract used in \hyperref[cp:panel]{the Taylor proof}. Input-span reduction
preserves the panel trajectory and all Gaussian first-row laws; it uses
no passive labels. The exact retained count remains

\[
1020(L+1)R^2+36R+10m(k+1)+pk+(m+p)+3(m+p)d.
\]

The additional gate

\begin{equation}
n\ge N_{\rm count}:=
 [32X^6(1+C_{\rm sel})\lambda^{-3/2}(1+Y^{-1})]^4
\label{cp:poly:eq:taylor-count-gate}
\end{equation}

implies $K+1\le8\ell_n$. To check it, the logarithm defining $K$ is
at most $(3/2)\log n+\sqrt{\ell_n}+\log(1+\sqrt{\ell_n})$ plus
the logarithm of the bracket in this gate. The latter is at most
$(1/4)\log n$, and the source gate gives $\ell_n\ge16$, so each
intervening term is at most $\ell_n/4$. Division by $\log2$ and the
ceiling leave the stated bound. The branches $\eta=Y$ and $\eta=1$
are smaller. Squaring the rank bound with
\eqref{cp:poly:eq:panel-count} and $k\le m+p$ yields

\[
R^2\le X^{810}(m+p)^2
 \left[(Y/\lambda)^4\ell_n^3+\lambda^{-2}\ell_n^2b_n^2\right].
\]

The exact inventory is therefore bounded by the Taylor headline with,
for example, coefficient $X^{820}$. The inverse-label gate belongs to
this simplification only. Without it the exact $K$ displays the extra
$\log(1/Y)$ dependence. Dense arrays, jets, coefficients and proof-panel
anchors are discarded after initialization, exactly as before.

We finally prove \Cref{cp:poly:width}. Allocate $\tau=\delta/2$
to the dense lower event and $\rho=\delta/4$ to the relevant source event.
On their intersection the Legendre ratio is at most

\[
4P^5a_\tau^{-1}\ell_n^{9/2}n^{-15/64}
 \le4P^5a_\tau^{-1}C(9/2,7/64)n^{-1/8}.
\]

Consequently the following fully specified maximum is a sufficient
Legendre width:

\begin{equation}
\left\lceil\max\left\{
 N_{\rm src}(m,\delta/4),N_{\rm low}(\delta/2),N_{\rm Leg},
 [4P^5C(9/2,7/64)/(a_{\delta/2}\varepsilon)]^8
 \right\}\right\rceil.
\label{cp:poly:eq:legendre-final-width}
\end{equation}

For Taylor the absolute error $Y/n$ and the same lower event give ratio
at most $a_\tau^{-1}\ell_n^{5/2}/\sqrt n\le64a_\tau^{-1}n^{-1/4}$.
A sufficient width is

\begin{equation}
\left\lceil\max\left\{
N_{\rm src}(m+p,\delta/4),N_{\rm low}(\delta/2),N_{\rm count},
[2C_{\rm sel}\lambda^{-1}(1+\lambda^{-1/2})]^{1/7},
[64/(a_{\delta/2}\varepsilon)]^4
\right\}\right\rceil.
\label{cp:poly:eq:taylor-final-width}
\end{equation}

The probability union costs less than $\delta$; no independence between
the source and comparison events is required. Since
$\lambda^{-1}=m/\gamma$, $u_\rho$ is logarithmic in $mL/\rho$,
$L\le X$, and $a_\tau^{-1}$ has only the displayed polynomial factors,
these maxima are bounded by the polynomial asserted in the theorem.
The proof counts retained real coordinates and exact-arithmetic error;
it does not infer efficient preprocessing or finite precision from a
polynomial sufficient width.
